\subsection{Well (WEL) Package}

The Well Package applies a specified volumetric flow rate to each listed cell.
The single per-boundary input is the well rate.

\begin{longtable}{p{0.15\textwidth} p{0.15\textwidth} p{0.60\textwidth}}
\caption{Memory-manager variables in the WEL Package that are candidates for
modification through the API. See also the variables common to all stress
packages described above.} \label{table:api-gwf-wel} \\

\hline
\textbf{Variable} & \textbf{Class} & \textbf{Guidance} \\
\hline
\endfirsthead

\hline
\textbf{Variable} & \textbf{Class} & \textbf{Guidance} \\
\hline
\endhead

\texttt{Q} & Period-reset &
Volumetric well rate for each well, dimensioned by \texttt{MAXBOUND} (positive
for injection, negative for extraction). Read from input at the start of each
stress period. Set it every time step (after \texttt{prepare\_time\_step}) to
change pumping at run time; a value set once before the time loop is overwritten
at the next Read and Prepare. \emph{Caveats:} when a well rate is supplied by a
time series, it is refreshed each time step and overwrites API changes (class
becomes Overwritten). When \texttt{AUTO\_FLOW\_REDUCE} is active, \texttt{Q} is
the requested rate; the rate actually applied to an extraction well is reduced as
the cell head approaches the cell bottom. \\

\texttt{AUXVAR} & Period-reset &
Auxiliary variable values, dimensioned by the number of auxiliary variables and
\texttt{MAXBOUND} (present when \texttt{AUXILIARY} variables are specified).
Re-read each stress period; set after \texttt{prepare\_time\_step}. Indexed by
boundary position, so a value follows the boundary index, not the cell. \\

\texttt{NODELIST} & Period-reset &
One-based reduced node number of the cell containing each boundary, dimensioned
by \texttt{MAXBOUND}. Set it after \texttt{prepare\_time\_step} to relocate a
boundary (see ``Changing Boundary Locations''). \\

\texttt{NBOUND} & Period-reset &
Number of active boundaries in the current stress period, from 0 to
\texttt{MAXBOUND}. Increase it (and fill the new entries) to activate boundaries
or decrease it to deactivate them (see ``Changing Boundary Locations''). \\

\hline
\end{longtable}
